\seccionreporte{Conclusiones y trabajo futuro}{sec:conclusiones}

\propositoseccion{%
  Cierra el hito: qué funciona hoy, qué falta para integrar con la app móvil y qué mejoras planean.%
}

\subsection{Conclusiones}

\begin{enumerate}
  \item \textbf{Logro principal:} Cumplimos el objetivo de automatizar el primer filtro de los proyectos. Ahora el sistema puede leer un documento y decirte matemáticamente si se parece a algo que ya existe, quitándole horas de trabajo a los profes.
  \item \textbf{Validación y persistencia:} El microservicio quedó súper sólido. Logramos meter las validaciones estrictas para rebotar archivos malos y todo queda guardado en la base de datos de SQLite para que haya trazabilidad de quién subió qué.
  \item \textbf{Lección aprendida:} Nos dimos cuenta de que no le puedes dar el control de las matemáticas a un LLM porque inventa números. Es mil veces mejor hacer el clustering duro con algoritmos como HDBSCAN y ya nada más usar el LLM para que redacte el reporte basándose en esos datos reales.
\end{enumerate}

\subsection{Integración con la aplicación móvil}

Para la siguiente fase, el flujo va a ser así: El alumno abre la app en Flutter y sube su PDF. La app hace una llamada HTTP POST a nuestro microservicio mandando el archivo. Cuando el backend regresa el JSON con el veredicto, la app lee el nivel de riesgo y pinta un semáforo en pantalla (Verde, Amarillo o Rojo) para avisarle al alumno si su proyecto está bien o si tiene que cambiarlo.

\subsection{Trabajo futuro}

Ahorita el MVP ya está funcional en Contabo, pero para dejarlo nivel empresarial nos falta:

\begin{itemize}
  \item Meterle balanceo de carga para que Ollama no truene si se conectan muchos al mismo tiempo.
  \item Migrar la base de datos de auditoría de SQLite a PostgreSQL cuando tengamos miles de registros.
  \item Ponerle HTTPS al servidor porque ahorita estamos pegándole por HTTP directo a la IP.
  \item Conectar todo esto finalmente con el frontend móvil en Flutter.
\end{itemize}

\subsection{Entregables finales del hito}

\begin{tabular}{@{}lp{8.5cm}@{}}
  \toprule
  Entregable & Estado \\
  \midrule
  Repositorio & \url{https://github.com/eduartrob/integratorProjectClustering-back-corvus} \\
  PDF (\LaTeX) & Este documento. \\
  Script Entrenamiento & \texttt{notebooks/run\_training.py} (en el repo). \\
  Diagrama arquitectura & PDF y PNG generados, en la carpeta correspondiente. \\
  API docs & Swagger interactivo en \texttt{/docs}. \\
  Colección pruebas & \texttt{collections/corvus\_clustering.postman\_collection.json}. \\
  Video 5--10 min & \ReporteVideo \\
  \bottomrule
\end{tabular}
